\scanpage{128}
\uline{Définition} Soit $K/\mathbb Q$ une extension \deleted{\unsure{alg. de gr.|de degré fini}} \inserted{algébrique}. On dit
que $K$ est tot. réel si pour tout $\mathbb Q$-hom $\sigma:K\to\mathbb C$,
on a $\sigma(K)\subset\mathbb R$. On dit qu'un élément $\unsure{a|x}$ de $K$ est
tot. réel \inserted{(resp. tot. $>0$, tot. $\ge0$)} si pour tt $\sigma$ comme dessus, on a
$\sigma(a)\in\mathbb R$ (resp. $\sigma(x)>0$, resp. $\sigma(a)\ge0$). Tot.\textsuperscript{t} (ou purement)
imaginaire : \ldots .

\uline{Sorite} Pour que $x$ soit tot. réel, il f. et s. que
$\mathbb Q(x)\subset K$ le soit. Pour qu'il soit tot. $\ge0$, il faut et
il suffit qu'il soit nul ou tot $>0$. Les éléments
tot $\ge0$ \deleted{\unsure{forment}} \inserted{sont} les éléments $\ge0$ pour une relation
d'ordre sur $K$. [On parlera de $x$ tot. $\ge y$ ou
$x$ tot. $\le y$]. Si $K$ est engendré par des ss-corps
$K_\alpha$, alors $K$ est tot. réel ssi les $K_\alpha$ le sont. Si
\inserted{$K$ est engendré par des $x_\alpha$,}
$K$ tot. réel ssi \deleted{\unsure{tous les}} les $x_\alpha$ le sont.

Il existe, pour $K$ donné, un plus grand ss-corps
totalement réel de $K$, appelé la \uline{clôture tot. réelle} de
$\mathbb Q$ dans $K$. Soit $\deleted{\unsure{K=|K}}\overline{\mathbb Q}\subset\mathbb C$, et $\pi=\operatorname{Gal}(\deleted{\unsure{K}}\overline{\mathbb Q}/\mathbb Q)$,
enfin $\tau\in\pi$ la conjugaison complexe
\[
\tau\lambda=\bar\lambda\qquad \text{pour }\lambda\in\overline{\mathbb Q}
\]
de sorte qu'on a
\[
\tau^2=1.
\]

Alors le ss-groupe $\pi'$ de $\pi$ qui \deleted{\unsure{invarie|invariant}} \inserted{correspond} à clôture
tot\textsuperscript{t} réelle $\deleted{\unsure{de|K}}\mathbb K\ \deleted{\unsure{\mathbb Q}}$ de $\mathbb Q$ dans $\overline{\mathbb Q}$ est \deleted{\unsure{le}} le ss-groupe
\margin{\diagram{\begin{tikzcd}[ampersand replacement=\&,row sep=small]A\arrow[d,dash]\\\unsure{L|\mathbb L}\arrow[d,dash]\\\unsure{K|\mathbb K}\arrow[d,dash]\\\mathbb Q\end{tikzcd}}}
\uline{invariant} \uline{fermé} engendré par $\tau$, i.e. le sous groupe invariant
engendré par les $ {}^{\sigma}\tau=\sigma\tau\sigma^{-1}$ $(\sigma\in\pi)$.

Considérons la plus grande extension abélienne
$A$ de $\mathbb K$, dont le groupe de Galois est donc $H\unsure{=|:}\pi'/(\pi',\pi')$. \deleted{\unsure{On|Son}} Sur
% archive sheet number: 127
